% Applied Glaciology -- Assignments part 1, shortest answers.
% Blue box: everything the sheet asks for that question. Answer: plain, in SI units.
% Grey box: the summarised answer in the units of the sheet.
% Numbers from ../code/si_units.csv and ../code/task*.csv; uncertainties (one standard deviation) and the
% significant figures they allow from ../code/errors.csv (code/errors.py). Build: latexmk min_si.tex
\documentclass{scrartcl}
\usepackage{agl}
% Builds from this one file, chosen by \aglmode: "full" (this file as it is, 23 pages), "report" (the title,
% contents, questions, references and figures: the pages before the supplement) and "supplement" (a cover, S1-S3
% and the closing views, with their own references); also "pruned" (the 15-page version: everything inside
% \cut{}, \swap{}{}, cutblock and cutfigure left out) and "marked" (the 23 pages with those passages struck
% through in red). min_si_report.tex, min_si_supplement.tex, min_si_15.tex and min_si_marked.tex set the mode.
\ifdefined\aglmode\else\def\aglmode{full}\fi
\usepackage{lua-ul}% \strikeThrough across paragraphs, lists and equations (not across images: see \agl@cross)
\makeatletter
\def\agl@modemarked{marked}\def\agl@modepruned{pruned}\def\agl@modereport{report}\def\agl@modesupplement{supplement}
\newif\ifagl@report\newif\ifagl@supplement
\ifx\aglmode\agl@modereport\agl@reporttrue\fi
\ifx\aglmode\agl@modesupplement\agl@supplementtrue\fi
\let\ifsupplementbuild\ifagl@supplement% for the body, outside \makeatletter
\newif\ifsplitbuild\ifagl@report\splitbuildtrue\fi\ifagl@supplement\splitbuildtrue\fi
% the two halves of the document; each split build drops the other half
\ifagl@supplement
  \NewDocumentEnvironment{reportpart}{+b}{}{}
\else
  \NewDocumentEnvironment{reportpart}{+b}{#1}{}
\fi
\ifagl@report
  \NewDocumentEnvironment{supplementpart}{+b}{}{}
  \usepackage{xr}\externaldocument[full-]{build/min_si}% the labels of the full build: Section S1 etc. still resolve
  \def\sref#1{full-#1}% a label of the other half
\else
  \NewDocumentEnvironment{supplementpart}{+b}{#1}{}
  \def\sref#1{#1}
\fi
\newcommand\agl@cross[1]{% a red cross over a box
  \begin{tikzpicture}[baseline=(i.base)]\node[inner sep=0pt] (i) {#1};
    \draw[accent, line width=1.2pt, opacity=.85] (i.south west) -- (i.north east);
    \draw[accent, line width=1.2pt, opacity=.85] (i.north west) -- (i.south east);
  \end{tikzpicture}}
\newcommand\agl@markgraphics{\ifdefined\agl@incg\else\let\agl@incg\includegraphics% once per nesting
  \renewcommand\includegraphics[2][]{\agl@cross{\agl@incg[##1]{##2}}}\fi}
\ifx\aglmode\agl@modemarked
  \NewDocumentCommand\cut{+m}{{\color{accent}\strikeThrough{#1}}}
  \NewDocumentCommand\swap{+m +m}{\cut{#1}{\color{accent}#2}}
  \NewDocumentCommand\variant{+m +m}{#1}
  \NewDocumentEnvironment{cutblock}{+b}{{\agl@markgraphics\cut{#1}}}{}
  \NewDocumentEnvironment{cutfigure}{O{tbp} +b}{\begin{figure}[#1]\agl@markgraphics\cut{#2}\end{figure}}{}
\else\ifx\aglmode\agl@modepruned
  \NewDocumentCommand\cut{+m}{}
  \NewDocumentCommand\swap{+m +m}{#2}
  \NewDocumentCommand\variant{+m +m}{#2}
  \renewcommand*\appendix{\agl@appendix\counterwithin{figure}{section}\counterwithin{table}{section}}% as agl.sty without its \clearpage: the supplements follow the references
  \NewDocumentEnvironment{cutblock}{+b}{}{}
  \NewDocumentEnvironment{cutfigure}{O{tbp} +b}{}{}
\else
  \NewDocumentCommand\cut{+m}{#1}
  \NewDocumentCommand\swap{+m +m}{#1}
  \NewDocumentCommand\variant{+m +m}{#1}
  \NewDocumentEnvironment{cutblock}{+b}{#1}{}
  \NewDocumentEnvironment{cutfigure}{O{tbp} +b}{\begin{figure}[#1]#2\end{figure}}{}
\fi\fi
\makeatother
\makeatletter% measure 128 mm, about 80 characters; figures keep 150 mm and overhang to the right
\agl@layout{portrait}{128mm}{242mm}{25mm}{8mm}{14mm}
\makeatother
\newcommand\wide[1]{\noindent\makebox[\textwidth][l]{#1}}
\addbibresource{refs.bib}
\ExecuteBibliographyOptions{giveninits=false}
\AtBeginDocument{\postdisplaypenalty=10000}% an equation stays with its symbols table
\setlength\emergencystretch{2.5em}
\usepackage{placeins}% \FloatBarrier
\renewcommand*\bottomfraction{.6}% the outline figure (0.54 of the page) may close the page of Question 5.3
\tcbuselibrary{breakable}% the question and unit boxes may break across pages, 3 lines at least before a break
\tcbset{every box/.style={breakable, lines before break=3}}
\RenewTColorBox{converted}{}{% grey: the summarised answer of each question in the units of the sheet
  colback=rule!20, colframe=rock, coltext=text,
  boxrule=0pt, leftrule=2pt, arc=0pt, outer arc=0pt, boxsep=0pt,
  left=3.6mm, right=3.6mm, top=2.2mm, bottom=2.2mm,
  before skip=.7\baselineskip, after skip=.7\baselineskip,
  before upper={\setlength\parindent{0pt}\setlength\parskip{.25\baselineskip}%
    {\sffamily\bfseries\color{rock}Summarised answer, in the units of the assignment}\par}}
\DeclareSIUnit\areal{\kilo\gram\per\square\metre}
\DeclareSIUnit\massflux{\kilo\gram\per\square\metre\per\second}
\DeclareSIUnit\ddfsi{\kilo\gram\per\square\metre\per\second\per\kelvin}
\groupmembers{Benjamin Fahl, Valentin Istomin}
\title{Ghiacciaio del Basòdino}
\subtitle{Mass balance, melt model and volume projection}
\date{21 October 2026}
\makeatletter% title page: the sky-view factor of the terrain fills the whole page (210 x 297 mm, one 0.5 m DEM pixel
% per image pixel); the title block, with the group and the date, stands on the glacier. The contents follow on page 2.
\newcommand*\agl@watermark[2]{%
  \node[anchor=south east, inner sep=2pt, fill=black, fill opacity=.25,
        text=white, text opacity=.8,
        font=\sffamily\fontsize{5}{6}\selectfont] at (#1.south east) {#2};}
\newcommand*\agl@bannerline{Sky-view factor of the terrain, swissALTI3D \qty{0.5}{\metre}, \textcopyright\,swisstopo; LV95 E 2678960--2681440, N 1139792--1143300}
\newcommand*\agl@titleblock{%
  \agl@label{Applied Glaciology\enskip\textperiodcentered\enskip Assignments part 1}\par
  \vspace{1.2mm}%
  \agl@label{Fall Semester 2026\enskip\textperiodcentered\enskip Submitted \@date}\par
  \vspace{3mm}%
  {\fontsize{44}{47}\selectfont\bfseries\color{ink}Ghiacciaio del\\Basòdino\par}%
  \vspace{2.5mm}%
  {\fontsize{14}{18}\fontseries{l}\selectfont\color{rock}\@subtitle\par}%
  \vspace{2.5mm}%
  {\color{ink}Benjamin Fahl and Valentin Istomin, in equal parts\enskip\textperiodcentered\enskip ETH Zürich\par}}
\renewcommand*\maketitle{%
  \thispagestyle{empty}\begingroup\setlength\parindent{0pt}\sffamily
  \begin{tikzpicture}[remember picture, overlay]
    \node[inner sep=0pt, anchor=north west] (p) at (current page.north west)
      {\includegraphics[width=\paperwidth]{figures/hero_page.jpg}};
    \agl@watermark{p}{\agl@bannerline}
    \node[anchor=north west, inner sep=0pt, text width=\textwidth, align=left]
      at ([xshift=\dimexpr 1in+\oddsidemargin\relax, yshift=-170mm]current page.north west) {\raggedright\agl@titleblock};
  \end{tikzpicture}%
  \null\endgroup\clearpage
  % contents, a page of their own: a hand-set list (topic, points, page), the orthophoto of the south margin at the foot
  \thispagestyle{empty}\begingroup\setlength\parindent{0pt}\sffamily\hypersetup{linkcolor=.}%
  \agl@label{Applied Glaciology\enskip\textperiodcentered\enskip Assignments part 1}\par
  \vspace{4mm}%
  {\fontsize{30}{34}\selectfont\bfseries\color{ink}Contents\par}%
  \vspace{7mm}%
  \agl@part{Questions}
  \agl@row{1}{Geodetic mass balance 1994--2021}{6 pts}{q:1}
  \agl@row{2}{The weather record since 1935}{5 pts}{q:2}
  \agl@row{3}{Degree-day model}{19 pts}{q:31}
  \agl@sub{3.1}{The model and its inputs}{7 pts}{q:31}
  \agl@sub{3.2}{Calibration, the factor and the literature}{6 pts}{q:32}
  \agl@sub{3.3}{Assumptions and their biases}{6 pts}{q:33}
  \agl@row{4}{Mass balance of every year 1935--2025}{7 pts}{q:4}
  \agl@row{5}{Volume and the years ahead}{14 pts}{q:51}
  \agl@sub{5.1}{Volume in 2021}{4 pts}{q:51}
  \agl@sub{5.2}{When half the ice is gone}{6 pts}{q:52}
  \agl@sub{5.3}{Beyond the constant area}{4 pts}{q:53}
  \agl@row{}{References}{}{sec:refs}
  \ifsplitbuild\agl@row{}{Our outline of 27 September 2026}{Figure 3}{fig:outline}\else
  \agl@row{}{The glacier from orbit\cut{, 1938--2026}}{Figure\swap{s 3--6}{ 3}}{fig:outline}\fi
  \vspace{3mm}%
  \agl@part{Supplements}
  \ifagl@report
  \agl@row{S1}{Area following the volume}{separate PDF}{}
  \agl@row{S2}{The bed from the thickness map}{separate PDF}{}
  \agl@row{S3}{Ice-flow model}{separate PDF}{}
  \agl@row{S4}{The glacier from orbit, 1938--2026}{separate PDF}{}
  \agl@row{}{The same view in 1850, 2026 and 2035}{separate PDF}{}
  \else
  \cut{\agl@row{S1}{Area following the volume}{}{sec:s1}}
  \agl@row{\swap{S2}{S1}}{The bed from the thickness map}{}{sec:s2}
  \agl@row{\swap{S3}{S2}}{Ice-flow model}{}{sec:s3}
  \cut{\agl@row{}{The same view in 1850, 2026 and 2035}{}{sec:outro}}
  \fi
  \begin{tikzpicture}[remember picture, overlay]
    \node[inner sep=0pt, anchor=south west] (b) at (current page.south west)
      {\includegraphics[width=\paperwidth]{figures/contents_banner.jpg}};
    \agl@watermark{b}{The south margin of the glacier, SWISSIMAGE \qty{10}{\centi\metre} of 2024 at \qty{0.2}{\metre} per pixel, \textcopyright\,swisstopo; LV95 E 2679804--2680796, N 1140292--1140708}
  \end{tikzpicture}%
  \endgroup\clearpage}
% rows of the contents page: number (blue), title, note (points) in grey, page number at the right
\newcommand*\agl@part[1]{%
  {\color{rule}\hrule height .4pt}\vspace{2.2mm}%
  \agl@label{#1}\par\vspace{2.8mm}}
\newcommand*\agl@row[4]{%
  \noindent\makebox[11mm][l]{\color{glacier}\bfseries\large #1}%
  {\color{ink}\bfseries\large #2}\hfill{\small\color{rock}#3}\hspace{5mm}\makebox[7mm][r]{\color{ink}\large\agl@pageof{#4}}\par
  \vspace{2.4mm}}
\newcommand*\agl@pageof[1]{\if\relax\detokenize{#1}\relax\else\pageref{#1}\fi}
\newcommand*\agl@sub[4]{%
  \noindent\hspace*{11mm}\makebox[11mm][l]{\color{glacier}#1}%
  {\color{ink}#2}\hfill{\small\color{rock}#3}\hspace{5mm}\makebox[7mm][r]{\color{ink}\pageref{#4}}\par
  \vspace{1.6mm}}
\ifagl@supplement% the supplement: a plain cover with the Robiei reservoir at the foot, then S1-S3 from page 2
\renewcommand*\maketitle{%
  \thispagestyle{empty}\begingroup\setlength\parindent{0pt}\sffamily
  \agl@label{Applied Glaciology\enskip\textperiodcentered\enskip Assignments part 1}\par
  \vspace{1.2mm}%
  \agl@label{Fall Semester 2026\enskip\textperiodcentered\enskip Submitted \@date}\par
  \vspace{3mm}%
  {\fontsize{44}{47}\selectfont\bfseries\color{ink}Ghiacciaio del\\Basòdino\par}%
  \vspace{2.5mm}%
  {\fontsize{14}{18}\fontseries{l}\selectfont\color{rock}Supplement to the report\par}%
  \vspace{2.5mm}%
  {\color{ink}Benjamin Fahl and Valentin Istomin, in equal parts\enskip\textperiodcentered\enskip ETH Zürich\par}%
  \vspace{12mm}%
  \agl@part{In this supplement}
  \agl@row{S1}{Area following the volume}{}{sec:s1}
  \agl@row{S2}{The bed from the thickness map}{}{sec:s2}
  \agl@row{S3}{Ice-flow model}{}{sec:s3}
  \agl@row{S4}{The glacier from orbit, 1938--2026}{}{sec:s4}
  \agl@row{}{References}{}{sec:srefs}
  \agl@row{}{The same view in 1850, 2026 and 2035}{}{sec:outro}
  \begin{tikzpicture}[remember picture, overlay]
    \node[inner sep=0pt, anchor=south west] (b) at (current page.south west)
      {\includegraphics[width=\paperwidth]{figures/supplement_banner.jpg}};
    \agl@watermark{b}{Sky-view factor of the Robiei reservoir and its dam, north of the glacier, swissALTI3D \qty{0.5}{\metre}, \textcopyright\,swisstopo; LV95 E 2680400--2682880, N 1144900--1145940}
  \end{tikzpicture}%
  \endgroup\clearpage}
\fi
\makeatother

\renewcommand\unv[1]{#1}% all passages accepted: set in black (remove this line to see them in blue again)
\newcommand\figsummers{% body of the figure: in the report of the unsplit builds, Section S4 of the supplement
  \wide{\includegraphics[width=143.5mm]{figures/satellite_summers.pdf}}% 96 %: at 150 mm the page is 9 mm short
  \credit{PlanetScope: \textcopyright\,Planet Labs PBC 2018--2026, Education and Research programme;
    Sentinel-2: contains modified Copernicus Sentinel data 2016--2024}
  \caption{The end of every summer 2016--2026 in true colour, PlanetScope (\qty{3}{\metre}) where
    the archive has a clear scene after 10 September and Sentinel-2 (\qty{10}{\metre}) for 2016,
    2017 and 2024. \unv{One colour stretch for all tiles; red: our outline of 27 September 2026 in
    every tile, so that the ice beyond it is the ice lost. The glacier of 2016 reaches well beyond
    the red line on every side; from 2022 on, the summers end with bare ice over the whole glacier.}}
  \label{fig:summers}
}
\newcommand\figrecordone{% body of the figure: in the report of the unsplit builds, Section S4 of the supplement
  \wide{\includegraphics[width=142.5mm]{figures/satellite_record_1.pdf}}% 95 %: at 150 mm the page is 10 mm short
  \credit{Aerial orthophotos: \textcopyright\,swisstopo (LUBIS, SWISSIMAGE); Landsat: U.S. Geological Survey;
    Sentinel-2: contains modified Copernicus Sentinel data}
  \caption{The end of every summer with a usable image, 1938--2015: the best image of each year after 10
    August, aerial orthophoto, Sentinel-2 or Landsat. \unv{The window E 2678800--2681800,
    N 1140200--1142200 holds the glacier of 1931. Red: our outline of 27 September 2026 in every tile.
    Black and white: aerial films; false colour (near infrared, red, green): the Landsat 5
    multispectral scanner at \qty{60}{\metre} in 2012; the other satellites in true colour with one
    stretch, near-saturated snow drawn grey. Years without a usable image are left out: no image, cloud,
    or a Landsat scene displaced by its geolocation (tier 2).}}
  \label{fig:record1}
}
\newcommand\figrecordtwo{% body of the figure: in the report of the unsplit builds, Section S4 of the supplement
  \wide{\includegraphics[width=135mm]{figures/satellite_record_2.pdf}}% 90 %
  \credit{Aerial orthophotos: \textcopyright\,swisstopo (SWISSIMAGE); Landsat: U.S. Geological Survey;
    Sentinel-2: contains modified Copernicus Sentinel data; PlanetScope: \textcopyright\,Planet Labs PBC}
  \caption{The end of every summer with a usable image, 2016--2026, continued from \cref{fig:record1}.}
  \label{fig:record2}
}
\begin{document}
\maketitle
\begin{reportpart}

% ---------------------------------------------------------------- 1
\begin{question}{1}{6}\label{q:1}
Use the geodetic method to compute the mass balance of Ghiacciaio del Basodino for
1994--2021, with a density of \qty{850}{\kilo\gram\per\cubic\metre} for the volume change.

\textbf{How large was the specific mass change rate during the period?} Express it in
m\,w.e.\,a$^{-1}$, and explain the assumptions you made.
\end{question}

\unv{We subtract the 1994 elevation model from the 2021 model, cell by cell, and add up the
change over every cell that was glacier in 1994 or 2021. Times the density this is the mass
change; divided by the mean glacier area and the length of the period it is the specific rate:
\[
  \dot B = \frac{\rho\,\Delta V}{\bar A\,\Delta t} ,
  \qquad
  \Delta V = a \sum_i \bigl(z_{2021,i} - z_{1994,i}\bigr) ,
  \qquad
  \bar A = \tfrac12 \bigl(A_{1994} + A_{2021}\bigr) .
\]
\begin{symbols}
  $\dot B$ & specific mass change rate 1994--2021 & \qty{-2.5 \pm 0.3 e-5}{\massflux} \\
  $\rho$ & density of the volume change (given) & \qty{850}{\kilo\gram\per\cubic\metre} \\
  $\Delta V$ & volume change, cells $i$ with ice in 1994 or 2021 & \qty{-5.0 \pm 0.5 e7}{\cubic\metre} \\
  $z_{1994,i}$, $z_{2021,i}$ & surface elevation of cell $i$ & \unit{\metre} \\
  $a$ & area of one cell, $\qty{10}{\metre} \times \qty{10}{\metre}$ & \qty{100}{\square\metre} \\
  $A_{1994}$, $A_{2021}$ & glacier area from the two masks, half a cell along the outline & \qty{2.40 \pm 0.05 e6}{}, \qty{1.61 \pm 0.05 e6}{\square\metre} \\
  $\bar A$ & mean glacier area (hint of the sheet) & \qty{2.00e6}{\square\metre} \\
  $\Delta t$ & October 1994 to October 2021, 27 years & \qty{8.52e8}{\second} \\
\end{symbols}
Over the whole period the glacier lost $\dot B\,\Delta t = \qty{-2.1 \pm 0.3 e4}{\areal}$.

Assumptions:
\begin{itemize}
  \item All of the volume change has the given density of \qty{850}{\kilo\gram\per\cubic\metre},
    which allows for some firn \parencite{Huss2013}.
  \item Both elevation models are from October, so the period holds 27 whole melt seasons.
  \item The two elevation models are aligned as supplied. Outside the glaciers they are
    identical except in 37 of \num{15775} cells at the ice edge, so the terrain cannot test the
    1994 heights over the glacier.
  \item The rate refers to the mean area. With the 1994 or the 2021 area alone it would be
    \qty{-2.1e-5}{} or \qty{-3.1e-5}{\massflux}.
\end{itemize}

Uncertainty, one standard deviation of each input: the 1994 model \qty{\pm 2}{\metre} and the
2021 model \qty{\pm 1}{\metre} as one offset over the glacier (our assumption; the sheet gives
no accuracy), \qty{11}{\percent} of the rate; the density \qty{\pm 60}{\kilo\gram\per\cubic\metre}
\parencite{Huss2013}, \qty{7}{\percent}; the masks half a cell along the outline,
\qty{2}{\percent}; the dates anywhere in October, \qty{0.1}{\percent}. Together
\qty{13}{\percent}. Glacier Monitoring Switzerland finds \qty{-2.2 \pm 0.7 e-5}{\massflux} for
1991--2021 from its own elevation models.

\begin{converted}
The glacier lost $0.78 \pm 0.10$ m\,w.e.\,a$^{-1}$ over 1994--2021, on the mean of the two
areas and at \qty{850}{\kilo\gram\per\cubic\metre}. A year has \qty{3.156e7}{\second}, and
\qty{1}{\kilo\gram} of water on \qty{1}{\square\metre} is a layer \qty{1}{\milli\metre} thick:
\begin{align*}
  \qty{-2.47e-5}{\massflux} \times \qty{3.156e7}{\second\per\annum}
  &= \qty{-779}{\kilo\gram\per\square\metre\per\annum} \\
  &= -0.78 \pm 0.10\ \text{m\,w.e.\,a}^{-1} .
\end{align*}
\begin{itemize}
  \item Whole period: $-21 \pm 3$ m\,w.e.
  \item With the 1994 area $-0.65$, with the 2021 area $-0.97$ m\,w.e.\,a$^{-1}$.
  \item Glacier Monitoring Switzerland, 1991--2021: $-0.69 \pm 0.21$ m\,w.e.\,a$^{-1}$.
\end{itemize}
\end{converted}}

% ---------------------------------------------------------------- 2
\begin{question}{2}{5}\label{q:2}
Retrieve a homogenized temperature and precipitation time series in monthly resolution, since
1935, for a meteorological station that you consider representative for the glacier.

\textbf{Which data did you choose and why?} State the url from which you downloaded the data,
and discuss the choice of the station.
\end{question}

\unv{We use the homogenised monthly temperature and precipitation of Locarno\mbox{-}Monti, a
MeteoSwiss station at \qty{367}{\metre}, \qty{36}{\kilo\metre} south-east of the glacier. The
series covers October 1934 to September 2025 without gaps. We downloaded it on 23 September 2026
from \url{https://data.geo.admin.ch/ch.meteoschweiz.ogd-nbcn/otl/ogd-nbcn_otl_m.csv}.

Three stations within \qty{40}{\kilo\metre} have homogenised temperature and precipitation
since 1935: Grimsel Hospiz, Andermatt and Locarno\mbox{-}Monti. Only Locarno\mbox{-}Monti lies south of the
main Alpine ridge, like the glacier, so it shares the glacier's precipitation from the south.
The drawback is height. The station lies \qtyrange{2.2}{2.8}{\kilo\metre} below the glacier,
so its temperature has to be carried up to every glacier cell with the lapse rate of Task~3,
\[
  T(z) = T_\text{station} - \Gamma\,(z - z_\text{station}) ,
  \qquad \Gamma = \qty{6.5e-3}{\kelvin\per\metre} ,
  \qquad z_\text{station} = \qty{367}{\metre} ,
\]
which makes the glacier \qtyrange{14}{18}{\kelvin} colder than the station. Any error in
$\Gamma$ grows with that height difference: \qty{0.5e-3}{\kelvin\per\metre} moves the glacier
temperature by \qty{1.3}{\kelvin} from Locarno\mbox{-}Monti and by \qty{0.5}{\kelvin} from Grimsel
Hospiz at \qty{1980}{\metre}. The measured balances of 1992--2025 cannot separate the three
stations: the model of Task~3 follows them with a correlation of 0.78 with Locarno\mbox{-}Monti,
0.74 with Grimsel Hospiz and 0.70 with Andermatt, and a correlation from 34 years is uncertain
by about $\pm 0.15$. The snow does: the winter precipitation of Locarno\mbox{-}Monti follows the
measured winter balance with a correlation of 0.67, that of Grimsel Hospiz, north of the ridge,
with 0.53.

\begin{converted}
Locarno\mbox{-}Monti, the MeteoSwiss station at \qty{367}{\metre}, \qty{36}{\kilo\metre}
south-east of the glacier and on its side of the Alpine ridge: homogenised monthly temperature
(\unit{\degreeCelsius}) and precipitation (\unit{\milli\metre}) from October 1934, downloaded on
23 September 2026 from the address above.
\end{converted}}

% ---------------------------------------------------------------- 3.1
\begin{question}{3.1}{7}\label{q:31}
Use the time series of Question~2 to calibrate the Degree Day Factor (DDF) of a simple
temperature index model.

\textbf{What was your procedure?} Give a step by step explanation and mention every assumption.
Discuss whether you use a distributed model (every grid cell separately) or not (for example
only the mean glacier elevation).
\end{question}

\unv{Melt $M$ is the mass of snow or ice that turns into water per square metre of glacier.
The model makes the melt rate proportional to the positive temperature $T^+$, the amount by
which the air is above the melting point, and zero below it:
\[
  T^+ = \max\bigl(T - T_\text{melt},\,0\bigr) , \qquad
  \frac{\mathrm{d}M}{\mathrm{d}t} = F\,T^+ .
\]
The factor $F$ is what we calibrate, with the monthly series of Question~2, October 1994 to
September 2021. Month by month and cell by cell, then year by year:
\begin{gather*}
  T_i(z) = T_{\text{station},i} - \Gamma\,(z - z_\text{station}) , \qquad
  T^+_i(z) = \max\bigl(T_i(z) - T_\text{melt},\,0\bigr) , \\
  P^\text{snow}_i(z) = \begin{cases} P_i & \text{if } T_i(z) < T_0 \\ 0 & \text{otherwise} \end{cases} , \qquad
  M_i(z) = F\,\tau\,T^+_i(z) , \\
  C_y = \Bigl\langle \sum_{i \in y} P^\text{snow}_i(z) \Bigr\rangle , \qquad
  S_y = \Bigl\langle \sum_{i \in y} \tau\,T^+_i(z) \Bigr\rangle , \qquad
  B_y = C_y - F\,S_y , \\
  F = \frac{\sum_y C_y - \dot B\,\Delta t}{\sum_y S_y} .
\end{gather*}
\begin{symbols}
  $T_i(z)$, $T^+_i(z)$ & temperature and its positive part & \unit{\kelvin} \\
  $T_{\text{station},i}$, $P_i$ & station temperature, precipitation & \unit{\kelvin}, \unit{\areal} \\
  $z$, $z_\text{station}$ & cell elevation (2021 surface); station & \qtyrange{2585}{3177}{}, \qty{367}{\metre} \\
  $\Gamma$ & lapse rate (given) & \qty{6.5e-3}{\kelvin\per\metre} \\
  $T_0$ & snow threshold (given) & \qty{275.15}{\kelvin} \\
  $T_\text{melt}$ & melt threshold (given) & \qty{273.15}{\kelvin} \\
  $\tau$ & month, $365/12$ days (given) & \qty{2.628e6}{\second} \\
  $P^\text{snow}_i(z)$, $M_i(z)$ & snowfall, melt of month $i$ at $z$ & \unit{\areal} \\
  $\langle\,\cdot\,\rangle$ & mean over the \num{16065} glacier cells & \\
  $C_y$ & snowfall of year $y$, glacier mean & \qty{1.23e3}{\areal}, mean year \\
  $S_y$ & positive temperature summed over year $y$ & \qty{4.52e7}{\kelvin\second}, mean year \\
  $B_y$ & balance of year $y$, glacier mean & \unit{\areal} \\
  $\dot B\,\Delta t$ & mass change 1994--2021 (Question~1) & \qty{-2.1 \pm 0.3 e4}{\areal} \\
  $F$ & degree-day factor & \qty{4.5 \pm 2.0 e-5}{\ddfsi} \\
\end{symbols}
The steps:
\begin{enumerate}
  \item Surface: the 2021 elevation model inside the 2021 mask, \num{16065} cells of
    $\qty{10}{\metre} \times \qty{10}{\metre}$, fixed in time; both models refer to October
    (given), so the period holds 27 whole balance years.
  \item For every cell $z$ and month $i$: the temperature from the station with the lapse rate;
    snow where the cell is colder than $T_0$, the same precipitation at every elevation (block
    rain model); melt from the positive temperature over the month (first two lines above). All
    given.
  \item For every hydrological year $y$: the glacier means of the snowfall, $C_y$, and of the
    positive temperature integrated over time, $S_y$; the balance $B_y = C_y - F S_y$.
  \item $F$ such that the 27 balances add up to the mass change of Question~1. $B_y$ is a
    straight line in $F$, so this gives $F$ directly (last line above).
\end{enumerate}
We compute every cell separately (distributed), not only the mean elevation, because snow and
melt switch on and off at thresholds and so do not change linearly with elevation. Here the
difference is small: at the mean elevation alone (\qty{2896}{\metre}) $F$ would be
\qty{1.3}{\percent} higher.

Assumptions beyond the given ones: the rate of Question~1, which refers to the mean area, is
applied to the 2021 surface (matching the total mass change on the 2021 area instead would
give $F$ \qty{9.5}{\percent} higher); meltwater does not refreeze; the station precipitation
needs no correction.

\begin{converted}
A distributed temperature-index model, every \qty{10}{\metre} cell of the 2021 surface, month
by month, with the given thresholds and lapse rate; $F$ is the one value for which the 27
modelled balances of October 1994 to September 2021 add up to the geodetic $-21 \pm 3$ m\,w.e.
of Question~1. Its value, in the unit of the sheet, is in Question~3.2.
\end{converted}}

% ---------------------------------------------------------------- 3.2
\begin{question}{3.2}{6}\label{q:32}
\textbf{What value (and unit) do you obtain for the calibrated DDF, and how does it compare to
values reported in the literature?}

Have \enquote{days} in the unit. Compare against at least two published DDF values, with
references, and discuss what might cause the differences.
\end{question}

\unv{$F = \qty{4.5 \pm 2.0 e-5}{\ddfsi}$; in the unit with days it is $3.8 \pm 1.7$ mm w.e. per
\unit{\degreeCelsius} per day (grey box). The uncertainty, with $F$ recalibrated for one
standard deviation of each input: lapse rate \qty{\pm 0.5e-3}{\kelvin\per\metre} (between
Locarno\mbox{-}Monti and Robiei, \qty{1898}{\metre}, the measured gradient of 1991--2025 is
\qtyrange{6.1}{7.4}{\kelvin\per\kilo\metre} in the melt months), \qty{\pm 1.8e-5}{};
precipitation \qty{\pm 25}{\percent}, \qty{\pm 0.7e-5}{}; the rate of Question~1
\qty{\pm 13}{\percent}, \qty{\pm 0.2e-5}{}; snow threshold \qty{\pm 1}{\kelvin},
\qty{\pm 0.2e-5}{\ddfsi}. The lapse rate dominates because it acts over \qty{2.5}{\kilo\metre}
of height.

Published values are \qty{6.1e-5}{} for snow and \qty{1.35e-4}{\ddfsi} for ice on
Aletschgletscher \parencite{Hock2003}, and \qty{4.1e-5}{} to \qty{5.3e-5}{\ddfsi} for snow on
66 glaciers \parencite{Braithwaite2008}. Ours agrees with the snow values within its
uncertainty and lies at a third of the ice value, although the glacier is bare ice for most of
the melt season: in 2026 the winter snow left it during July, and from August on it melted as
ice (\cref{fig:year}).

$F$ is fitted, so it takes up the errors of the input:
\begin{itemize}
  \item The temperature comes from a station \qty{2.5}{\kilo\metre} lower. The glacier gets
    many warm months, and more warm months need a smaller $F$ for the same melt. With Grimsel
    Hospiz, \qty{0.9}{\kilo\metre} below the glacier but north of the ridge, $F$ would be
    \qty{6.7e-5}{\ddfsi}; its snow fits the glacier worst (Question~2).
  \item The model gets only \qty{60}{\percent} of the winter snow measured on the glacier
    (\qty{880}{} against \qty{1470}{\areal}). Less snow needs less melt, so $F$ is smaller.
  \item One $F$ serves snow and ice, while the published values separate them.
  \item The published values come from daily data at single points, ours from monthly means
    over the whole glacier.
\end{itemize}

\begin{converted}
DDF $= 3.8 \pm 1.7$ mm\,w.e.\,(\unit{\degreeCelsius}\,d)$^{-1}$, within the published snow
values and a third of the ice value. A day has \qty{86400}{\second}:
\[
  F = \qty{4.46e-5}{\ddfsi} \times \qty{86400}{\second\per\day}
  = 3.85\ \text{mm\,w.e.}\,(\unit{\degreeCelsius}\,\text{d})^{-1} ,
\]
that is $3.8 \pm 1.7$ at one standard deviation.
The published values with the same factor, in mm\,w.e.\,(\unit{\degreeCelsius}\,d)$^{-1}$:
\begin{itemize}
  \item Aletschgletscher: 5.3 (snow), 11.7 (ice).
  \item 66 glaciers: 3.5 to 4.6 (snow).
  \item Our model with Grimsel Hospiz: 5.8.
\end{itemize}
\end{converted}}

% Figure 1 follows the answer that refers to it. The grey box above ends its page exactly, so the
% float would only be registered on the next page and its page would split Question 4: the barrier
% makes it the next page. (If the grey box ever stops filling its page, the barrier ends that page
% early; then drop it.)
\begin{figure}[p]
  \wide{\includegraphics[width=148.5mm]{figures/satellite_year_2026_planet.pdf}}% 99 %: at 150 mm the page is 1.2 mm short
  \credit{PlanetScope: \textcopyright\,Planet Labs PBC 2025--2026, Education and Research programme;
    Sentinel-2: contains modified Copernicus Sentinel data 2026, surface reflectance (Level-2A) from
    Earth Search, Element 84}
  \caption{Ghiacciaio del Basòdino through the hydrological year 2026 in true colour, PlanetScope
    (\qty{3}{\metre}) where a clear scene exists and Sentinel-2 (\qty{10}{\metre}) otherwise. \unv{All
    tiles with one colour stretch; red: our glacier outline of 27 September 2026. Fresh snow on 8
    October 2025; the glacier still white on 23 May and 13 June; bare ice spreading from the north
    during July; bare ice everywhere from August; on 12 October 2026, the newest clear scene, the
    first snow of the new hydrological year covers the glacier again. The coloured stripe on 13 June
    is a sensor artefact.}}
  \label{fig:year}
\end{figure}

% ---------------------------------------------------------------- 3.3
\begin{question}{3.3}{6}\label{q:33}
\textbf{Reflect upon the assumptions that were necessary for the calibration. Which are likely
to introduce a bias? In what direction do the individual biases go, and which assumptions are
likely to introduce the largest biases?}
\end{question}

\unv{Largest first:
\begin{itemize}
  \item Station temperature carried up \qty{2.5}{\kilo\metre} with one fixed lapse rate. The
    glacier is probably too warm, so $F$ is too low. The choice of station alone moves $F$
    between \qty{4.46e-5}{} and \qty{8.09e-5}{\ddfsi}; a lapse rate
    \qty{0.5e-3}{\kelvin\per\metre} off moves it by \qty{40}{\percent}.
  \item The same precipitation at every elevation, without the snow that the gauge misses or
    that wind and avalanches bring. The model has too little snow, so $F$ is too low.
  \item Monthly means hide warm days in cold months. Daily temperatures give
    \qty{10}{\percent} more degree-days, so $F$ is about \qty{10}{\percent} too high.
  \item The surface is fixed at 2021, although the glacier reached lower before. The model
    counts too little warmth, so $F$ is \qtyrange{3}{5}{\percent} too high.
  \item Density, \qty{\pm 60}{\kilo\gram\per\cubic\metre}: $F$ changes by
    \qty{\pm 3}{\percent}, either way.
  \item One factor for snow and ice. This does not shift $F$, but the model melts too much in
    snowy years and too little in bare-ice years, so its extreme years are too mild.
\end{itemize}
The two largest biases come from the input data, temperature and snow, and both push $F$
down.

\begin{converted}
The largest biases, the temperature carried up \qty{2.5}{\kilo\metre} and the block-rain snow,
both make $F$ too low. The station range, in mm\,w.e.\,(\unit{\degreeCelsius}\,d)$^{-1}$:
\begin{itemize}
  \item $\qty{4.46e-5}{\ddfsi} \times \qty{86400}{\second\per\day} = 3.85$,
  \item $\qty{8.09e-5}{\ddfsi} \times \qty{86400}{\second\per\day} = 6.99$.
\end{itemize}
\end{converted}}

% ---------------------------------------------------------------- 4
\begin{question}{4}{7}\label{q:4}
Use the model of Task~3 and the same time series to reconstruct the annual mass balance for
1935--2025, in m\,w.e.\,a$^{-1}$, for hydrological years (1~Oct.\ -- 30~Sep.). Plot the
results.

\textbf{What are (i) the cumulative and (ii) the average mass balance, and what do you observe
in the yearly series?} Can you discern trends, and do periods or single years stand out?
\end{question}

\begin{figure}[bp]% placed before the answer; at the foot of a page, so that the unbreakable start of Question 5.1 keeps its room
  \wide{\includegraphics{figures/mb_1935_2025_si.pdf}}
  \caption{Mass balance of each hydrological year, 1935--2025; \qty{1e3}{\areal} is
    \qty{1}{\metre} of water. \unv{(a) Modelled (bars) with its 11-year running mean, and
    measured since 1992 (circles) and over the survey periods (lines). (b) Cumulative balance,
    modelled and measured.}}
  \label{fig:mb}
\end{figure}

\unv{We run the model of Question~3.1 with the calibrated $F$ on the fixed 2021 surface for the
91 hydrological years 1935--2025 (\cref{fig:mb}).

\noindent\makebox[2.4em][l]{(i)}Cumulative: \qty{-1.0 \pm 2.0 e4}{\areal}. \\
\makebox[2.4em][l]{(ii)}Average: \qty{-0.4 \pm 0.7 e-5}{\massflux}.

The uncertainty comes from the inputs of Question~3 with $F$ recalibrated each time (lapse
rate \qty{\pm 1.5e4}{}, precipitation \qty{\pm 1.1e4}{}, the rate of Question~1
\qty{\pm 0.7e4}{}, snow threshold \qty{\pm 0.5e4}{\areal}): the calibration fixes the 27 years
1995--2021, so the error sits in the other 64. The geodetic surveys of Glacier Monitoring
Switzerland, chained with its measured years, give \qty{-4.8 \pm 0.7 e4}{\areal} for
1935--2025. Before 1962 the model is too positive by \qty{1.1e3}{\areal} per year; the glacier of that time was
larger and reached lower than the 2021 surface the model runs on (Hint~2).

\begin{itemize}
  \item The glacier gained mass until the 1980s and has lost it since. The mean rate was
    \qty{+1.2 \pm 1.3 e-5}{\massflux} in 1935--1965, \qty{+0.5 \pm 0.9 e-5}{} in 1966--1995 and
    \qty{-2.8 \pm 0.4 e-5}{} in 1996--2025. The cumulative balance peaks in 1984. Since 1996
    only two years are positive.
  \item Trend: none until 1985; from 1986 the rate falls by \num{3.3 \pm 1.1 e-14}
    \unit{\kilo\gram\per\square\metre\per\square\second} (\num{\pm 0.7e-14} from the fit
    alone); the grey box gives it in the unit of the sheet.
  \item Single years: the most negative are 2022 (\qty{-2.3e3}{\areal}), 2003 and 2018, all
    hot summers; 2018 and 2022 also followed dry winters. The most positive is 1977
    (\qty{+2.3e3}{\areal}), the wettest year of the series.
  \item Check: against the 34 balances measured on the glacier since 1992 the correlation is
    0.78 (0.59 to 0.88 at \qty{95}{\percent}) and the root mean square difference
    \qty{0.55e3}{\areal} per year, but the extremes are too mild. In 2022 the measured loss was
    \qty{3.5e3}{\areal}.
\end{itemize}

\begin{converted}
\begin{itemize}
  \item (i) $\qty{-1.02e4}{\areal} \div 1000 = -10.2$, that is $-10 \pm 20$ m\,w.e.; the
    measured chain $-48 \pm 7$ m\,w.e.
  \item (ii) $\qty{-3.57e-6}{\massflux} \times \qty{3.156e7}{\second\per\annum}
    = \qty{-113}{\kilo\gram\per\square\metre\per\annum}$,\\ that is $-0.11 \pm 0.22$ m\,w.e.\,a$^{-1}$.
  \item Period means $+0.38 \pm 0.40$, $+0.17 \pm 0.27$ and $-0.90 \pm 0.11$ m\,w.e.\,a$^{-1}$.
  \item Trend from 1986 on: times \qty{3.156e8}{\second} per decade, then as a rate:
    $-0.33 \pm 0.11$ m\,w.e.\,a$^{-1}$ per decade.
  \item Years: 2022 $-2.3$ (measured $-3.5$), 1977 $+2.3$ m\,w.e.; a modelled year is
    uncertain by $0.55$ m\,w.e.
\end{itemize}
\end{converted}}


% ---------------------------------------------------------------- 5.1
\begin{question}{5.1}{4}\label{q:51}
Use volume-area scaling to estimate the volume of the glacier for the year 2021.

\textbf{What was the glacier volume in 2021?} Express the result in \unit{\kmcubed}.
\end{question}

\unv{Volume--area scaling with the 2021 area of Question~1:
\[
  V_{2021} = c\,A_{2021}^{\gamma} = 0.19\,\sqrt[4]{\mathrm{m}} \times
  \bigl(\qty{1.61e6}{\square\metre}\bigr)^{1.375} = \qty{6.5e7}{\cubic\metre} .
\]
\begin{symbols}
  $V_{2021}$ & glacier volume in 2021 & \qty{6.5 \pm 2.2 e7}{\cubic\metre} \\
  $A_{2021}$ & glacier area in 2021 (Question~1) & \qty{1.61 \pm 0.05 e6}{\square\metre} \\
  $c$ & scaling constant, worldwide mean \parencite{Bahr2015}, unit $\mathrm{m}^{3 - 2\gamma}$ & $0.19\,\sqrt[4]{\mathrm{m}}$, $\pm \qty{34}{\percent}$ \\
  $\gamma$ & scaling exponent, fixed by theory \parencite{Bahr1997} & 1.375 \\
\end{symbols}
The mean thickness $V/A = c\,A^{0.375}$ is then \qty{41 \pm 14}{\metre}. The relation gives the
average glacier of this area, not this glacier: one standard deviation of $c$ changes the
volume by \qty{34}{\percent} \parencite{Bahr2015}, half a cell along the outline by
\qty{4.7}{\percent}, so $V_{2021}$ is uncertain by \qty{34}{\percent}. The measured thickness
map of Swiss glaciers \parencite{Grab2021} gives only \qty{1.4e7}{\cubic\metre} for 2021
(Question~5.3).

\begin{converted}
$V_{2021} = 0.065 \pm 0.022$ km$^3$, a mean thickness of $41 \pm 14$ m. A cubic kilometre is
\qty{1e9}{\cubic\metre}.
\begin{itemize}
  \item $V_{2021} = \qty{6.5e7}{\cubic\metre} \times 10^{-9} = \qty{0.065 \pm 0.022}{\kmcubed}$.
  \item $c$ in the length unit of $A$: $0.19\,\mathrm{m}^{1/4} = 0.034\,\mathrm{km}^{1/4}$.
\end{itemize}
\end{converted}}

% ---------------------------------------------------------------- 5.2
\begin{question}{5.2}{6}\label{q:52}
Use the mass balance series of Task~4 to compute the mass change until 2025. For the future,
assume that the mass balance rate stays at its average of 1995--2025. Assume that the glacier
area stays at its 2021 value (Hint~1).

\textbf{In which year will the glacier have lost \qty{50}{\percent} of its 2021 volume?} Discuss
whether this projection is likely an under- or overestimate, and why.
\end{question}

\unv{Mass per area becomes ice thickness by the ice density, and volume by the constant area:
\[
  V(t) = V_{2021} + \frac{A_{2021}}{\rho_\text{ice}}
    \Bigl(\sum_{y=2022}^{2025} B_y + (t - t_{2025})\,\dot B_\text{future}\Bigr) ,
  \qquad
  V(t_{1/2}) = \tfrac12\,V_{2021} .
\]
\begin{symbols}
  $V(t)$ & glacier volume at time $t$ & \unit{\cubic\metre} \\
  $V_{2021}$ & volume in 2021 (Question~5.1) & \qty{6.5 \pm 2.2 e7}{\cubic\metre} \\
  $A_{2021}$ & area in 2021, kept constant (Hint~1) & \qty{1.61e6}{\square\metre} \\
  $\rho_\text{ice}$ & mean glacier density (Hint~4) & \qty{900}{\kilo\gram\per\cubic\metre} \\
  $\sum B_y$ & balances of 2022--2025 from Task~4 & \qty{-5.9 \pm 1.1 e3}{\areal} \\
  $\dot B_\text{future}$ & mean rate of the 30 years 1996--2025 & \qty{-2.8 \pm 0.4 e-5}{\massflux} \\
  $t_{2025}$ & end of September 2025 & \\
  $t_{1/2}$ & time at which half of $V_{2021}$ is lost & $2039 \pm 8$ \\
\end{symbols}
Until 2025 the glacier loses \qty{1.0e7}{\cubic\metre}, which leaves
\qty{5.5e7}{\cubic\metre}; after that it loses \qty{1.6e6}{\cubic\metre} per year. Half of the
2021 volume is gone in 2039, uncertain by 8 years (one standard deviation): 7 from the
scaling volume, 2 from the rate of Question~1 through the future balance, 1.5 from the lapse
rate and the precipitation. Read as the 31 years 1995--2025, the mean rate is
\qty{-2.75e-5}{\massflux} and the year 2040, a difference far inside that uncertainty.

We think 2039 is too late, so the projection overestimates the time left:
\begin{itemize}
  \item Constant area (Hint~1) makes the year too early. A shrinking glacier loses less ice
    each year, and it loses its lowest, most negative parts first. If the area follows the
    volume through the scaling relation, the year is 2044.
  \item Constant balance makes the year too late, because the climate keeps warming. With the
    mean rate of 2015--2025, \qty{-4.1e-5}{\massflux}, the year is 2035.
  \item Task~4 loses too little in 2022--2025: \qty{5.9e3}{} against \qty{8.6e3}{\areal}
    measured. With the measured values the year is 2036.
  \item The scaling volume is probably too large (Question~5.1). Half of the
    \qty{1.4e7}{\cubic\metre} from the thickness map is less than the ice lost by 2025.
\end{itemize}
Three of the four effects make the year earlier. Each shift is smaller than the uncertainty of
the year itself, which the 2021 volume sets.

\begin{converted}
Half of the 2021 volume is gone in $2039 \pm 8$; the projection overestimates the time left.
The future rate:
\begin{align*}
  \dot B_\text{future} &= \qty{-2.84e-5}{\massflux} \times \qty{3.156e7}{\second\per\annum}
  \div 1000 \\
  &= -0.90 \pm 0.11\ \text{m\,w.e.\,a}^{-1} .
\end{align*}
\begin{itemize}
  \item Over 31 years $-0.87$; for 2015--2025 $-1.3$ m\,w.e.\,a$^{-1}$.
  \item 2022--2025: $-5.9 \pm 1.1$ m\,w.e. modelled, $-8.6$ measured.
  \item Volumes times $10^{-9}$: $0.065 \pm 0.022$ in 2021, 0.055 in 2025, half of 2021
    \qty{0.033}{\kmcubed}.
  \item Loss after 2025: \qty{0.0016}{\kmcubed} per year.
\end{itemize}
\end{converted}}

% ---------------------------------------------------------------- 5.3
\begin{question}{5.3}{4}\label{q:53}
\textbf{Can you think of a strategy that avoids the constant area of Hint~1?}

Explain how, and what additional information, tools or methods (if any) you would need.
\end{question}

\unv{Three ways, from simple to complete. We tried all three; the supplement shows each.
\begin{enumerate}
  \item Let the area follow the volume. Every year, compute the area from the volume with the
    scaling relation of Question~5.1, $A = (V/c)^{1/\gamma}$. This needs nothing new. Half of the
    2021 volume is then gone in 2044 instead of 2039\cut{ (\cref{\sref{sec:s1}})}.
  \item Use the bed. The ice thickness of 2021 is the surface minus the bed. Every year each
    cell loses the balance of its elevation band, and a cell leaves the glacier when its ice is
    gone. This needs the bed, which the thickness map of all Swiss glaciers \parencite{Grab2021}
    provides, for Basòdino from radar profiles of 2006. Melted with the balances measured since
    2021, the map loses area faster than the satellite shows, so its ice is too thin: the
    observed retreat fits a 2021 volume of \qty{2.1e7}{} to \qty{2.4e7}{\cubic\metre}
    (\cref{\sref{sec:s2}}). By the measured balances the glacier has lost
    \qty{1.70e7}{\cubic\metre} since September 2021, 1.2 times the volume the map gave it, and
    it is still there.
  \item An ice-flow model, which also moves ice from the top towards the tongue. This needs the
    bed, a flow law and far more computing. Our run set of \num{10062} ice-flow projections puts
    the glacier below \qty{1e4}{\square\metre} in 2039 to 2040, the same in all three emission
    scenarios, and in \qty{96}{\percent} of the runs in the same year with and without flow
    (\cref{\sref{sec:s3}}). For a glacier \qty{9}{\metre} thick the flow no longer matters; the bed
    does.
\end{enumerate}
\ifsplitbuild Satellite images check the area every year (supplement, Section~S4). \else\cut{Satellite images check the area every year (\cref{fig:summers,fig:record1,fig:record2}). }\fi With Sentinel-2 we measured
\qty{1.03 \pm 0.09 e6}{\square\metre} on 27 September 2026 (\cref{fig:outline}),
\qty{35 \pm 8}{\percent} less than in 2021 by the same method.

\begin{converted}
Area from volume by scaling (half gone in 2044); the measured bed melted cell by cell; an
ice-flow model (ice gone in 2039 to 2040); the last two need the bed and computing.
\begin{itemize}
  \item Areas times $10^{-6}$: \qty{1.03e6}{\square\metre} is \qty{1.03}{\kilo\metre\squared};
    \qty{1e4}{\square\metre} is \qty{0.01}{\kilo\metre\squared}.
  \item Volumes times $10^{-9}$: \qty{1.4e7}{\cubic\metre} is \qty{0.014}{\kmcubed};
    \qty{1.70e7}{\cubic\metre} lost is \qty{0.017}{\kmcubed}.
  \item \qty{2.1e7}{} to \qty{2.4e7}{\cubic\metre} is \qty{0.021}{} to \qty{0.024}{\kmcubed}.
\end{itemize}
\end{converted}}

% The figures of Question 5.3 follow its answer and are numbered in this order: the outline closes
% the page on which the answer ends, the two tall record figures fill the next two pages, record 2
% closes the references page.
\begin{figure}[bp]% foot of the page on which Question 5.3 ends
  \wide{\includegraphics[width=144mm]{figures/outline_2026_si.pdf}}% 96 %: at 150 mm the foot of that page is 3 mm short
  \credit{Contains modified Copernicus Sentinel data 2026; surface reflectance (Level-2A) from
    Earth Search, Element 84; 2024 outline: swissTLM3D, \textcopyright\,swisstopo; 2021 outline:
    course material}
  \caption{Our outline of Ghiacciaio del Basòdino on 27 September 2026, over Sentinel-2 true
    colour of that day. \unv{Red: \qty{1.03 \pm 0.09 e6}{\square\metre}, a body of
    \qty{1.02e6}{\square\metre} with four holes of rock or debris and two detached blocks of
    ice; the uncertainty is one \qty{10}{\metre} pixel along the whole perimeter. Dark blue: the
    2024 outline of the Federal Office of Topography (\qty{1.38e6}{\square\metre}); dashed: the
    course mask of 2021 (\qty{1.61e6}{\square\metre}).}}
  \label{fig:outline}
\end{figure}

\ifsplitbuild\else% the three satellite figures stay in the report only when it is not split
\begin{cutfigure}[p]\figsummers\end{cutfigure}
\begin{cutfigure}[p]\figrecordone\end{cutfigure}
\begin{cutfigure}[bp]\figrecordtwo\end{cutfigure}% foot of the references page; at 150 mm it would be half a line too tall for it
\fi

\phantomsection\label{sec:refs}
\printbibliography[title=References]
\end{reportpart}

% ------------------------------------------------------- supplement
\begin{supplementpart}
\appendix
\renewcommand*\thesection{S\arabic{section}}
\renewcommand*\bottomfraction{.75}% a tall figure may close the last text page of a section
\begin{cutblock}
\section{Area following the volume}\label{sec:s1}

\unv{After 2025 the volume of each hydrological year is the previous one minus the ice removed
by the balance of Question~5.2, $A\,\dot B / \rho_\text{ice}$, with the area of that year; the
next area is $A = (V/c)^{1/\gamma}$ with the $c$ and $\gamma$ of Question~5.1\cut{
(\cref{fig:s1})}. Half of the 2021 volume is gone in 2044. In 2100 \qty{3.4e6}{\cubic\metre}
remain, where the constant area has no ice after 2060. Both curves keep the balance rate fixed
while the glacier retreats; in reality the lowest parts, which lose the most, go first, and the
balance of the remaining glacier is less negative.\cut{\par Computed with \texttt{code/task5\_projection.py}, Python on a laptop, under one second.}}

\begin{cutfigure}[htbp]
  \wide{\includegraphics{figures/volume_projection_si.pdf}}
  \caption{Projected volume at the end of each hydrological year. \unv{Scaling for 2021, the
    balances of 2022--2025 from Task~4 (shaded), then the mean balance of 1996--2025. Black:
    area constant at 2021; blue: area following the volume. Dashed: half of the 2021 volume;
    dots: the first year below it.}}
  \label{fig:s1}
\end{cutfigure}
\end{cutblock}

\section{The bed from the thickness map}\label{sec:s2}

\unv{The Swiss thickness map \parencite{Grab2021, Grab2020data} gives the bed and the ice
thickness on a \qty{10}{\metre} grid for the surface of 2012, for Basòdino from six radar
profiles of 2006\cut{ (\cref{fig:s2})}. The thickness on 30 September 2021 is the 2021 surface minus
the bed, zero where the bed lies above the surface (\qty{13}{\percent} of the 2021 area):
\qty{1.4e7}{\cubic\metre}, \qty{9}{\metre} on average. The map's own uncertainty is wide: cell
by cell it spans \qty{0.6e7}{} to \qty{4.3e7}{\cubic\metre} around its 2012 volume of
\qty{2.3e7}{\cubic\metre}.

\cut{From there we melted every cell with the balance that Glacier Monitoring Switzerland measured
in its \qty{100}{\metre} elevation band, period by period, from September 2021 to September
2026, times $0.86 \pm 0.12$ (the ratio of the geodetic loss to the sum of the stake balances
over 2012--2024). A cell leaves the glacier when its ice is gone. The ice area left is compared with
the glacier outlines we drew from Sentinel-2:
\par\medskip\noindent{\small\addfontfeature{Numbers=Tabular}
\begin{tabular}{@{} l S[table-format=1.2] S[table-format=1.2] @{}}\toprule
  ice area (\qty{e6}{\square\metre}) & {map, as published} & {observed} \\
  \midrule
  30 September 2021 & 1.38 & 1.58 \\
  September 2024    & 1.03 & 1.39 \\
  September 2025    & 0.92 & 1.24 \\
  September 2026    & 0.60 & 1.03 \\
  \bottomrule
\end{tabular}}\par\medskip}%
The map loses area too fast, so its ice is too thin. Thickened everywhere by 0.4 to 0.6 of its
stated uncertainty, it matches the outlines of 2025 and 2026, and the 2021 volume becomes
\qty{2.1e7}{} to \qty{2.4e7}{\cubic\metre}: 1.5 to 1.7 times the map, a third of the scaling
value. Beyond 2026 this method needs modelled balances per band; the runs without flow of
\cref{sec:s3} are exactly that.

\cut{The volume alone shows the same. From September 2021 to August 2024 Glacier Monitoring
Switzerland measured $-1.99 \pm 0.25$ m\,w.e. per year (geodetic), then \qty{-1.52}{} m\,w.e.
in 2025 and \qty{-3.50}{} m\,w.e. in 2026 (provisional). Over the 2021 area these balances remove
\qty{1.70 \pm 0.16 e7}{\cubic\metre} of ice between 30 September 2021 and 29 September 2026:
$1.20 \pm 0.11$ times the volume the map gave the glacier in 2021 (the uncertainty is that of
the loss; the map's own is what the test is about). By the map the glacier
would have vanished during the summer of 2026; on 27 September it covered
\qty{1.03e6}{\square\metre}. The same test on every Swiss glacier with a thickness map and a
later survey, 1313 of 1400: by the latest survey, 2019 to 2024, 18 had lost more ice than the
map holds, two of them beyond two standard deviations; Basòdino stood at $0.91 \pm 0.13$ in
August 2024, before the summers of 2025 and 2026.\par Computed with \texttt{thickness.py}, \texttt{area\_clock.py}, \texttt{volume\_clock.py} and
\texttt{national\_screen.py} in \texttt{code/}, Python on a laptop, seconds.}}

\begin{cutfigure}[htbp]
  \wide{\includegraphics{figures/thickness_sections_si.pdf}}
  \credit{Ice thickness and bed: \textcite{Grab2020data}, Creative Commons Attribution 4.0;
    elevation models and masks: course material}
  \caption{Published ice thickness for the surface of 2012, with two sections. \unv{(a)
    Thickness over the outline of the Swiss Glacier Inventory 2016; black: the 2021 outline.
    (b) Section from the summit to the tongue, true scale. (c) Section across the main body
    along two radar profiles of 2006, vertical scale three times the horizontal. Grey: rock;
    pale blue: ice of 2021; dashed: the 1994 surface.}}
  \label{fig:s2}
\end{cutfigure}

\section{Ice-flow model}\label{sec:s3}

\unv{We used the Instructed Glacier Model, version 3.2.0 \parencite{Jouvet2023}, on the
\qty{10}{\metre} grid of the course data, \num{340} by \num{290} cells. Its mass balance is a
temperature-index model with a radiation term, driven by the daily weather of Locarno\mbox{-}Monti
carried up to the glacier with lapse rates measured against the Robiei station below the
glacier, with snow moved by the curvature of the surface. The balance parameters were
calibrated on the measurements of Glacier Monitoring Switzerland 1991--2026\cut{ (winter and
annual balances per elevation band, 344 stake readings, and the geodetic volume changes)}:
\num{131072} parameter draws, resampled to 128 by their likelihood.

Each run starts on 30 September 2021 from the measured thickness of \cref{sec:s2}, thickened
by a drawn share of its uncertainty (2021 volume \qty{1.9e7}{} to \qty{2.5e7}{\cubic\metre}),
runs through the observed weather to September 2026, and then under one of 26 climate model
chains of the Swiss scenarios CH2025 for one of three emission scenarios (SSP1-2.6, SSP2-4.5,
SSP5-8.5) to 2100: 128 draws by 26 chains by 3 scenarios, \num{9984} members\cut{, plus 78 with the
map as published}.\cut{ Each chain is shifted so that its 2026 matches the observed station
temperature.} Every member was also run without ice flow.

Results (\swap{\cref{fig:s3flow,fig:s3a,fig:s3b,fig:s3c}}{\cref{fig:s3a}}):
\begin{itemize}
  \item Half of the 2021 volume is gone in 2026, in every member: the observed weather of
    2021--2026 did that. Half of the 2026 volume is gone in 2029 (2028 to 2032 for \qty{90}{\percent}
    of the members).
  \item The ice area falls below \qty{1e4}{\square\metre} in 2039 to 2040 (median), 2035 to
    2048 for \qty{90}{\percent} of the members. The three emission scenarios differ by at most
    1.5 years, because the ice is gone before they separate.\cut{ With the chains left unshifted,
    about \qty{0.5}{\kelvin} cooler in 2026, the median is 2047.}
  \item With and without flow, \qty{96}{\percent} of the members disappear in the same year.
\cut{    The spread comes from the climate chain, the 2026 temperature level and the order of the
    weather years, not from the flow. The modelled glacier of 1850 flows at \qty{17}{\metre}
    per year in the median and up to \qty{27}{\metre} per year in its central basin; the
    hindcast's own ice of 2021 at \qty{1.7}{\metre} per year; the ice restarted from the measured
    thickness at \qty{0.25}{\metre} per year in 2021 and \qty{0.04}{\metre} per year by 2026
    (\cref{fig:s3flow}).}
  \cut{\item Check against the past: the same model, run from 1850, loses \qty{1.4e7}{\cubic\metre}
    between 1907 and 1991, against \qty{1.6 \pm 0.4 e8}{\cubic\metre} measured (six surveys
    chained), and hits only two of
    the four inventory areas. The projections start from the measured 2021 state and do not use
    this run, but the model does not reproduce the glacier of the twentieth century.}
\end{itemize}
\cut{\par Computed on the server s03 with two NVIDIA H200 graphics cards, eight runs per card at once:
the \num{10062} projections from 18:11 on 9 October to 02:18 on 10 October 2026, 8.1 hours,
97 process-hours, a median of \qty{32}{\second} per run; the \num{9984} runs without flow in 8
minutes. Code in \texttt{igm/}.}}

\begin{cutfigure}[tbp]
  \wide{\includegraphics[width=138mm]{figures/igm_flow.pdf}}% 92 %: with the text of S3 on one page
  \caption{Flow of the modelled glacier, the film member of the run set. \unv{(a) Surface speed
    of the 1850 state in steps of \qty{5}{\metre} per year with streamlines of the velocity, over
    a hillshade of the bed; red: the 1850 outline of the inventory; dashed: section A. (b) Ice
    thickness along section A, summit to tongue: the hindcast ice lost between 1850, 1900, 2000
    and 2021 as grey bands, the hindcast's own 2021 ice in blue, the 2021 start from the measured
    thickness in red; the upper axis gives the bed elevation. (c) Surface speed over the
    ice-covered cells, weighted by thickness, median and \qtyrange{5}{95}{\percent}: the hindcast
    1850--2021 in blue, the restart from the measured thickness 2021--2038 in red.}}
  \label{fig:s3flow}
\end{cutfigure}

\newcommand\figsthreea{% its own page; with the text in the 15-page version
  \variant{\wide{\includegraphics{figures/igm_volume_area.pdf}}}{\wide{\includegraphics[width=110mm]{figures/igm_volume_area.pdf}}}
  \caption{Volume and ice area 2021--2100 per emission scenario. \unv{Blue: median and
    \qtyrange{17}{83}{\percent} and \qtyrange{5}{95}{\percent} of the members with the chains
    shifted to the observed 2026; grey: median and \qty{95}{\percent} with the chains unshifted.
    Black: the 2021 start and the observed areas 2024--2026. The time axis is compressed after
    2050.}}
  \label{fig:s3a}
}
\variant{\begin{figure}[p]\figsthreea\end{figure}}{\begin{figure}[htbp]\figsthreea\end{figure}}

\begin{cutfigure}[p]
  \wide{\includegraphics{figures/igm_icefree_map.pdf}}
  \caption{Year each cell becomes ice-free, chains shifted to 2026, all scenarios. \unv{(a)
    Median over the members; (b) the range of \qty{90}{\percent} of the members. Outlines of
    the glacier inventories 1850 to 2024 and ours of 2026.}}
  \label{fig:s3b}
\end{cutfigure}

\newcommand\figsthreec{% in the supplement at the top of a page, the satellite section below it
  \wide{\includegraphics{figures/igm_years.pdf}}
  \caption{The years of the ice-flow run set. \unv{(a) Half of the 2026 volume, area below
    \qty{1e5}{} and below \qty{1e4}{\square\metre}, per scenario: 5, 17, 50, 83 and
    95 per cent of the members; blue: chains shifted to 2026, grey: unshifted. (b) Share of the
    spread per factor. (c) Share of the members whose year changes without flow.}}
  \label{fig:s3c}
}
\ifsupplementbuild\begin{figure}[tbp]\figsthreec\end{figure}\else
\begin{cutfigure}[p]\figsthreec\end{cutfigure}\fi

\ifsupplementbuild
\FloatBarrier
\section{The glacier from orbit, 1938--2026}\label{sec:s4}
The end of every summer with a usable image, 1938 to 2026, with our outline of 27 September 2026 in every tile
(\cref{fig:summers,fig:record1,fig:record2}; the outline itself is Figure~3 of the report). With Sentinel-2 we
measured \qty{1.03 \pm 0.09 e6}{\square\metre} on that day, \qty{35 \pm 8}{\percent} less than in 2021 by the same method.
\begin{figure}[p]\figsummers\end{figure}
\begin{figure}[p]\figrecordone\end{figure}
\begin{figure}[p]\figrecordtwo\end{figure}
\FloatBarrier\phantomsection\label{sec:srefs}\printbibliography[title=References]\fi
% ------------------------------------------------------------- outro
\begin{cutblock}
\clearpage
\thispagestyle{bare}\KOMAoptions{headsepline=false}
\vspace*{\fill}\phantomsection\label{sec:outro}
\wide{\includegraphics[width=150mm]{figures/outro_1850.jpg}}\par
\vspace{2.5mm}
\wide{\includegraphics[width=150mm]{figures/outro_2026.jpg}}\par
\vspace{2.5mm}
\wide{\includegraphics[width=150mm]{figures/outro_2035.jpg}}\par
\credit{Ice-flow model (\cref{\sref{sec:s3}}); terrain: swissALTI3D and SWISSIMAGE, \textcopyright\,swisstopo; Blender}
{\sffamily\small\color{rock}\noindent\unv{The same camera from the south-east in 1850, 2026 and 2035. Top: the
modelled glacier of 1850, \qty{2.6e8}{\cubic\metre} of ice, run from a steady state through the observed weather.
Middle: the model's 2026 state, restarted from the measured thickness in 2021. Bottom: 2035, under the middle
emission scenario, four years before the ice area falls below \qty{1e4}{\square\metre}. The smooth grey surface
is the modelled bed where the thickness map covers, not terrain.}\par}
\vspace*{\fill}
\end{cutblock}
\end{supplementpart}

\end{document}
